\section{Evaluation}

The evaluation tests ASIEP as a local evidence-conformance layer, not as a benchmark of agent capability. We do not evaluate whether ASIEP improves an agent's task performance, planning ability, or reasoning quality. Instead, we evaluate whether the current minimal implementation can make agent self-improvement evidence checkable: distinguishing valid and invalid evidence records, resolving local evidence bundles, detecting known tampering cases, generating bounded repair plans, packaging validated bundles, and reporting cross-standard mapping coverage. This framing is intentionally narrow. The evaluation measures auditability, closure, and conformance behavior on a curated local fixture corpus. It does not claim external certification, production deployment, or general benchmark performance.

The evaluation corpus contains valid ASIEP evidence records, invalid evidence records, local evidence bundles, missing-artifact bundles, digest-mismatch bundles, path-escape bundles, OpenTelemetry-like and LangSmith-like import fixtures, package requests, and adversarial evidence cases. The evaluation pipeline reuses the implemented tools. The validator checks profile conformance, lifecycle consistency, evidence references, gate consistency, rollback evidence, and semantic invariants. The repairer produces bounded repair plans for invalid records. The resolver verifies local artifact bindings and digest consistency. The importer generates ASIEP evidence bundles from local trace-like fixtures under a reference-only policy. The packager creates local FDO-like records, RO-Crate-like metadata, PROV JSON-LD, and package manifests. The evaluator aggregates these results into machine-readable reports and paper-ready tables.

The reported metrics are evidence-infrastructure metrics rather than task-performance metrics.

\begin{table*}[t]
\centering
\caption{Local fixture evaluation metrics. AUTHOR\_LAYOUT\_CHECK\_REQUIRED}
\begin{tabular}{p{0.25\linewidth}p{0.08\linewidth}p{0.58\linewidth}}
\toprule
Metric & Value & Interpretation \\
\midrule
\texttt{evidence\_completeness} & 1.0 & Required fields and artifact roles are present in known valid local cases. \\
\texttt{cross\_standard\_coverage} & 0.9333 & Most ASIEP field groups have local mappings to PROV, OpenTelemetry-like, LangSmith-like, FDO-like, RO-Crate-like, and governance categories. \\
\texttt{gate\_reproducibility} & 1.0 & Gate decisions in the local corpus can be checked from gate reports, thresholds, flip counts, and safety checks. \\
\texttt{tamper\_detection\_recall} & 1.0 & Known invalid and adversarial samples are detected by the implemented local checks. \\
\texttt{false\_positive\_rate} & 0.0 & Known valid local fixtures are not rejected by the toolchain. \\
\texttt{privacy\_policy\_compliance} & 1.0 & Generated valid packages avoid default embedding of raw prompts, user inputs, and model outputs under the current sentinel/key scanning policy. \\
\texttt{packaging\_closure} & 1.0 & Generated packages contain the expected manifest, evidence record, bundle, artifacts, FDO-like record, RO-Crate-like metadata, and PROV JSON-LD. \\
\texttt{agent\_readability} & 1.0 & Profile manifests, validator outputs, repair plans, bundle resolution results, import results, package results, and evaluation reports are machine-readable. \\
\bottomrule
\end{tabular}
\end{table*}


The evidence paths for these results are \texttt{reports/asiep\_v0.1\_evaluation\_report.json}, \texttt{paper\_assets/tables/evaluation\_metrics.md}, \texttt{paper\_assets/tables/attack\_corpus\_results.md}, \texttt{evaluation/corpus/asiep\_v0.1\_corpus.json}, and \texttt{evaluation/crosswalk/asiep\_v0.1\_crosswalk.json}.

On the current local corpus, all known valid examples pass the relevant checks, and all known invalid or adversarial examples are detected by at least one component of the toolchain. Missing gate reports, unsafe promotion decisions, pass-to-fail threshold violations, broken hash chains, missing artifacts, digest mismatches, and path-escape attempts are surfaced as structured errors. The validator reports these failures through stable error codes, JSON paths, invariant identifiers, repairability classes, and remediation hints. This allows the evaluator to treat validation as a reproducible machine-readable process rather than as a manual inspection step.

The generated repair plans preserve the evidence boundary. They request missing evidence, recommend lifecycle corrections when supporting evidence already exists, or suggest safe decision changes when a promotion is inconsistent with the gate result. They do not fabricate gate reports, rewrite safety results, change a safety regression from \texttt{true} to \texttt{false}, or forge digests to match tampered artifacts. This distinction is important because a repair mechanism for evidence objects should improve structural consistency without weakening evidential integrity.

The resolver results show that ASIEP references can be converted into locally verifiable bindings. A reference is not accepted merely because it appears in the profile. The resolver checks that the artifact exists inside the local bundle, remains within the bundle boundary, and matches the declared digest. The packager results further show that validated and resolved bundles can be converted into local research-object outputs, including package manifests, FDO-like records, RO-Crate-like metadata, and PROV JSON-LD.

These results indicate that ASIEP can support reproducible local evidence checking for agent self-improvement events. In particular, the combination of lifecycle invariants, digest-bound bundle resolution, and gate-consistency checks provides a practical way to distinguish valid records from several common evidence failures. The evaluation also shows that ASIEP can be imported from local trace-like fixtures and repackaged as local research objects.

The evaluation remains limited. The corpus is local and curated, and the reported values should not be interpreted as general benchmark results. The OpenTelemetry-like and LangSmith-like traces are local fixtures rather than live integrations. The FDO-like and RO-Crate-like packages are local representations rather than externally certified records. Privacy-policy compliance is checked through sentinel and key scanning, not full data-loss prevention. These limitations are intentional for v0.1: the goal is to demonstrate a reproducible evidence profile and local toolchain before claiming external deployment, third-party assurance, or broad coverage across agent systems.
